\SourceLeaf{246}%
% A544-P0246-D01
Num congruentia proposita divisores aequales habeat, eodem modo diiudicatur, uti in aequationum theoria. Ponamus
\[
X\equiv\xi^mP
\]
patet fore
\[
\frac{\mathrm{d}X}{\mathrm{d}x}\equiv\xi^{m-1}\left(mP\frac{\mathrm{d}\xi}{\mathrm{d}x}+\xi\frac{\mathrm{d}P}{\mathrm{d}x}\right)
\]
quare $\frac{\mathrm{d}X}{\mathrm{d}x}$ per $\xi^{m-1}$ dividetur. Generaliter sit
\[
X\equiv A^aB^bC^c\text{ etc.}
\]
ubi $A,B,C$ etc. denotant functiones primas diversas, erit
\[
\frac{\mathrm{d}X}{\mathrm{d}x}\equiv X\left(\frac{a\,\mathrm{d}A}{A\,\mathrm{d}x}+\frac{b\,\mathrm{d}B}{B\,\mathrm{d}x}+\frac{c\,\mathrm{d}C}{C\,\mathrm{d}x}+\text{etc.}\right)
\]
unde patet, nisi aliquis numerorum $a,b,c$ etc. per modulum dividatur, $\frac{\mathrm{d}X}{\mathrm{d}x}$ per $A^{a-1}B^{b-1}C^{c-1}$ etc. dividi posse, non autem per $A^a,B^b,C^c$ etc. Hinc sequitur

% A544-P0246-D02
\textsc{Theorema.} \emph{Si functionum $X$ et $\frac{\mathrm{d}X}{\mathrm{d}x}$ divisor communis maximae dimensionis sit $\xi$, omnes factores primos, quos $\xi$ habet, etiam $X$ habebit et quidem quemvis toties $+1$ vice quoties $\xi$, si igitur $X$ et $\frac{\mathrm{d}X}{\mathrm{d}x}$ sint functiones inter se primae, $X$ nullos factores aequales habebit.}

% A544-P0246-D03
\Article{369}
Exemplum I. Quaeritur an functio
\[
x^5+3x^4-6x^3+3x-4\ldots(X)
\]
secundum modulum $17$ divisores aequales habeat. Erit
\[
\frac{\mathrm{d}X}{\mathrm{d}x}\equiv5x^4-5x^3-xx+3
\]
Hinc invenitur, functiones $X$ et $\frac{\mathrm{d}X}{\mathrm{d}x}$ inter se esse primas, quare $X$ divisores aequales non habet.

% A544-P0246-D04
Exemplum II. Sit
\[
X\equiv x^5+6x^4-3x^3-4xx+2x-3(\mathrm{mod.}\,13)
\]
erit
\[
\frac{\mathrm{d}X}{\mathrm{d}x}\equiv5x^4-2x^3+4xx+5x+2
\]
maxima vero functionum $X,\frac{\mathrm{d}X}{\mathrm{d}x}$ communis mensura $\equiv5xx+7x+7$ seu 
\SourceLeaf{247}%
% A544-P0247-D01
multiplicata per $8$: $xx+4x+4$; at quum hic divisor sit $\equiv(x+2)^2$, functio $X$ per $(x+2)^3$ dividi poterit quotiensque (qui est $xx+11$) nullum amplius divisorem duplicem involvit.

% A544-P0247-D02
\Article{370.\ 371}
Si ex \textsection.\textsection.\ praecc. functio $X$ ita est exhibita $A^aB^bC^c$ etc., ita ut $A,B,C$ etc. inter se sint primae et numeri $a,b,c$ etc. inaequales, resolutio etiam ulterius extendi potest. Sit itaque $X$ functio, quae nullos amplius divisores aequales involvit. Supra vidimus, $x^p-x$ esse productum ex omnibus functionibus primis unius dimensionis. Sit $\xi$ divisor communis maximae dimensionis functionum $X$ et $x^p-x$, erit $\xi$ productum ex omnibus divisoribus ipsius $X$ unius dimensionis, et $\frac X\xi$ huiusmodi divisores non amplius habebit. Quodsi autem inveniatur, functiones $X$ et $x^p-x$ esse inter se primas, $X$ nullum divisorem unius dimensionis habebit adeoque congruentia $X\equiv0$ radices reales non habebit. Porro quoniam $x^{pp}-x$ est productum ex omnibus functionibus primis duarum dimensionum uniusque, divisor communis maximae dimensionis functionum $x^{pp}-x$ et $\frac X\xi$, $\xi'$ involvet omnes divisores ipsius $X$, qui sunt duarum dimensionum. Hinc ulterius progrediendo perspicitur, $X$ hoc modo in factores $\xi,\xi',\xi''$ etc. resolvi, qui continent respective omnes divisores unius, duarum, trium etc. dimensionum.

% A544-P0247-D03
\Article{372}
Si autem productum ex pluribus functionibus primis eiusdem dimensionis datum est, singulae functiones tentando erui debent. Magnam analogiam habet hoc problema cum eo, quod numerorum compositorum factores quaerere iubet. Hic vero iam a priori determinatur, an functio proposita in factores adhuc discerpi possit. Quum et hic factorum omnium possibilium multitudo sit finita, simili subsidio ut supra uti possumus. Sed huic rei inhaerere nolumus, nam calculator exercitatus principia probe assecutus, quando opus est, facile artificia particularia reperiet.

% A544-P0247-D04
Progredimur ad aliud caput, scilicet ad considerationem congruentiarum, si modulus non est numerus primus, uti hactenus semper supposuimus. Praesertim vero hic ille casus attentione dignus est, ubi modulus est numeri primi potestas, tum per se tum quod ad aliqua dubia removenda (\textsection.\textsection.\ .\ .) necessarius sit.

\SourceLeaf{248}%
% A544-P0248-D01
\Article{373}
\textsc{Problema.} \emph{Si functio $X$ secundum modulum $p$ in factores inter se primos $\xi,\xi',\xi''$ etc. sit resoluta, $X$ secundum modulum $pp$ in similes factores $\Xi,\Xi',\Xi''$ etc. resolvere ita, ut sit}
\[
\xi\equiv\Xi,\quad\xi'\equiv\Xi',\quad\xi''\equiv\Xi'',\text{ etc.}\;(\mathrm{mod.}\,p)
\]

% A544-P0248-D02
\emph{Sol.} Sit $X\equiv\xi\psi(\mathrm{mod.}\,p)$ seu $X=\xi\psi+p\Sigma$. Ponatur
\[
\Xi=\xi+p\varphi,\quad\Psi=\psi+p\omega
\]
erit
\[
\Xi\Psi=X-p\Sigma+(\varphi\psi+\xi\omega)p+pp\varphi\omega
\]
Si igitur $\Xi\Psi$ esse debet $\equiv X(\mathrm{mod.}\,pp)$, necessario debet esse $\varphi\psi+\xi\omega-\Sigma$ per $p$ divisibilis. At cum $\psi$ et $\xi$ secundum modulum $p$ sint functiones inter se primae, $\varphi$ et $\omega$ ita determinari poterunt, ut haec conditio adimpleatur (\textsection.\ 336), et quidem insuper ita, ut dimensiones ipsarum $\varphi$ et $\omega$ sint respective unitate minores dimensionibus functionum $\xi,\psi$. Hinc erit $X\equiv\Xi\Psi(\mathrm{mod.}\,pp)$. Patet, simili modo $\Psi$ rursus in factores $\Xi'\Omega$ discerpi posse, ita ut alter $\Xi'$ sit $\equiv\xi'(\mathrm{mod.}\,p)$ et ita porro, unde tandem
\[
X\equiv\Xi\,\Xi'\,\Xi''\text{ etc.}\;(\mathrm{mod.}\,pp).\quad\text{Q. E. Fac.}
\]

% A544-P0248-D03
\Article{374}
Facile hinc probari potest, functionem $X$ etiam secundum modulos $p^3,p^4$ etc. in factores resolvi posse. Generaliter sit
\[
X\equiv PQ(\mathrm{mod.}\,p^m)\quad\text{seu}\quad X=PQ+p^mR
\]
et functio $P$ ad ipsam $Q$ prima secundum modulum $p$; posito
\[
P'=P+Ap^m,\quad Q'=Q+Bp^m
\]
erit
\[
P'Q'=X-p^mR+(AQ+BP)p^m+ABp^{2m}
\]
Hinc pro quovis modulo $p^\nu$ ($\nu$ existente $>m$ et $<2m+1$) erit
\[
P'Q'\equiv X,\quad\text{si}\quad R\equiv AQ+BP(\mathrm{mod.}\,p^{\nu-m})
\]

\SourceLeaf{249}%
% A544-P0249-D01
Ex his perspicitur, si functio $X$ aequales non habeat divisores secundum modulum $p$, eam secundum modulum $p^k$ similiter in factores discerpi posse, uti secundum modulum $p$. At si $X$ divisores aequales habeat, res fit multo magis complicata neque adeo ex principiis praecedentibus prorsus exhauriri potest. Quare quum quae huc pertineant cuncta communicare non possimus, unicum casum tantummodo considerabimus, qui plurimum occurrit cuiusque enodatio ad quaedam in praecedentibus dubia solvenda requiritur. Hic est, si factores aequales unius dimensionis tantum respiciantur. Hic proprie etiam ad congruentiarum radices inveniendas adhiberi potest. Generaliter alia occasione hanc rem pertractabimus.

% A544-P0249-D02
\Article{375}
Sit igitur $X\equiv X'(x-a)^m(\mathrm{mod.}\,p)$ et functio $X'$ ad $x-a$ prima: desiderantur omnes divisores unius dimensionis huic $x-a$ secundum modulum $p$ congrui ipsius $X$ secundum modulos $pp,p^3$ etc. (Supponimus, functionem $X$ absolute per $x-a$ dividi non posse; alias enim $x-a$ secundum modulum quemcunque functionem $X$ divideret). Si substituatur $z+a$ pro $x$, habebitur
\[
Z\equiv Z'z^m(\mathrm{mod.}\,p)\quad\text{seu}\quad Z=Z'z^m+pA
\]
Iam si $Z$ secundum modulum $pp$ per aliquem divisorem formae $z+\alpha p$ dividi potest, necessario $A$ debet esse formae $zZ''+pB$. Nisi hoc sit, disquisitio iam est finita. Ponamus igitur
\[
\begin{aligned}
Z&\equiv Z'z^m+pZ''z(\mathrm{mod.}\,pp)\\
&\text{seu}\quad Z=Z'z^m+pZ''z+ppB
\end{aligned}
\]
patetque, $Z$ per $z$ ac quemcunque alium divisorem huic secundum modulum $p$ congruum dividi posse.

% A544-P0249-D03
Ut attentio fixetur, ponemus $m=4$, facile perspicitur, quemvis alium casum simili modo tractari posse. Iam si $Z$ secundum modulum $p^3$ per aliquem divisorem formae $z+\alpha p$ dividi potest, erit
\[
\begin{aligned}
0&\equiv-\alpha ppZ''+ppB(\mathrm{modd.}\,z+\alpha p,p^3)\\
&\text{seu}\quad\alpha Z''\equiv B(\mathrm{modd.}\,z,p)
\end{aligned}
\]
Iam tres casus esse possunt

% A544-P0249-D04
1) si $Z''\equiv0(\mathrm{modd.}\,z,p)$ et $B$ non $\equiv0$, tunc patet, nullum ipsius $\alpha$ valorem congruentiae satisfacere adeoque $Z$ secundum modulum $p^3$ nullum divisorem formae $z+\alpha p$ habere. Quare disquisitio erit finita

\SourceLeaf{250}%
% A544-P0250-D01
2) si nec $Z''$ nec $B\equiv0(\mathrm{modd.}\,z,p)$; tunc $\alpha$ unicum valorem habebit, scilicet
\[
\alpha\equiv\frac B{Z''}(\mathrm{modd.}\,z,p).
\]
Quare erit unicus divisor $\equiv z+\alpha p(\mathrm{mod.}\,pp)$ ipsius $Z$ secundum modulum $p^3$; eritque
\[
Z\equiv V(z+\alpha p)+p^3W
\]
Iam ponatur divisor ipsius $Z(\mathrm{mod.}\,p^4)$ $z+\alpha p+\beta pp$ eritque
\[
0\equiv\qquad\qquad
\]
\BookRule{.28\linewidth}

% A544-P0250-D02
\par\bigskip\Needspace{6\baselineskip}
\pdfbookmark[0]{Bemerkungen zur Analysis residuorum}{bemerkungen-analysis-residuorum}
\begin{center}\large\bfseries BEMERKUNGEN ZUR ANALYSIS RESIDUORUM.\end{center}
\BookRule{.09\linewidth}
\renewcommand{\BookRunningTitle}{Nachlass. Bemerkungen zur Analysis residuorum}

% A544-P0250-D03
Die beiden vorstehenden Abhandlungen sind einem umfangreichen Manu\-scripte entnommen, welches den Titel Analysis Residuorum führt und vermuthlich aus dem Jahre 1797 oder 1798 stammt; durch eine gänzliche Umarbeitung sind aus demselben später die Disquisitiones Arithmeticae entstanden. Der vollständige Titel des Caput sextum lautet:

% A544-P0250-D04
Solutio congruentiae $x^m-1\equiv0$ et aequationis $x^m-1=0$; cum dilucidationibus super theoria polygonorum regularium.

% A544-P0250-D05
Der zweite Theil desselben (\textsection\textsection.\ 253--278) ist seinem wesentlichen Inhalte nach in die siebente Section der Disqq. Arithm. übergegangen.

% A544-P0250-D06
Ausserdem ist noch zum Theil erhalten das Caput septimum. Variae quarundam investigationum praecedentium applicationes (\textsection\textsection.\ 279--302). Es zerfällt in folgende Unterabtheilungen:

% A544-P0250-D07
De fractionum communium transmutationibus (\textsection\textsection.\ 279--281).

% A544-P0250-D08
De fractionum communium in decimales conversione (\textsection\textsection.\ 282--292).

% A544-P0250-D09
De resolutione aequationis indeterminatae $xx=a+by$ (\textsection\textsection.\ 293--297).

% A544-P0250-D10
De resolutione aequationis indeterminatae $axx+byy=c$ (\textsection\textsection.\ 298--301).

% A544-P0250-D11
De investigatione divisorum numerorum (\textsection.\ 302; die folgenden Bogen fehlen).

% A544-P0250-D12
Dies alles ist fast wörtlich in die sechste Section der Disqq. Arithm. aufgenommen.

% A544-P0250-D13
Die beiden hier mitgetheilten Abschnitte behandeln die Gegenstände, welche, wie aus der Vorrede und den Artikeln 11, 44, 61, 62, 65, 84 der Disqq. Arithm. hervorgeht, den Inhalt der achten Section dieses Werkes bilden sollten. Es verdient indessen bemerkt zu werden, dass dieser Plan später wieder abgeändert ist; es findet sich nemlich unter den Manuscripten ein Fragment mit der Ueberschrift Sectio octava: Quarundam disquisitionum ad circuli sectionem pertinentium uberior consideratio. Dasselbe 
